% ----------------------------------------------------------------------------------------------------------------------------------
% docDynamo PDF template
% ----------------------------------------------------------------------------------------------------------------------------------
\documentclass[letterpaper,12pt]{article}

% Add hyperlinks to TOC
% ----------------------------------------------------------------------------------------------------------------------------------
\usepackage{hyperref}

% Allow EPS files
% ----------------------------------------------------------------------------------------------------------------------------------
\usepackage{graphicx}
\usepackage{epstopdf}

% Use the caption package to enable captions that are not numbered (caption*)
% ----------------------------------------------------------------------------------------------------------------------------------
\usepackage[font=small,textfont=it,justification=justified,singlelinecheck=false]{caption}

% Create a light gray color to use for source code listings
% ----------------------------------------------------------------------------------------------------------------------------------
\usepackage[table]{xcolor}
\definecolor{ltgray}{HTML}{E8E8E8}
\definecolor{dkblue}{HTML}{396A93}

% Styling for hyperlinks
% ----------------------------------------------------------------------------------------------------------------------------------
\hypersetup{frenchlinks=true}
% {
%    colorlinks,
%    linkcolor={dkblue},
%    citecolor={dkblue},
%    urlcolor={dkblue}
%}

% Use listings package instead of verbatim for displaying code
% ----------------------------------------------------------------------------------------------------------------------------------
\usepackage{courier}

\usepackage[warn]{textcomp}
\usepackage{listings}
\lstset
{
    basicstyle=\small\ttfamily,
    columns=flexible,
    breaklines=true,
    frame=tb,
    backgroundcolor=\color{ltgray},
    upquote=true
}

% Use tabularx for tables
% ----------------------------------------------------------------------------------------------------------------------------------
\usepackage{tabularx}
% \usepackage{ltablex}  This cannot be used because there is no package for RHEL/CENTOS
\newcolumntype{R}{>{\raggedleft\arraybackslash}X}%
\renewcommand{\arraystretch}{1.3}

% Allow four section levels (The fourth is implemented with paragraph)
% ----------------------------------------------------------------------------------------------------------------------------------
\usepackage{titlesec}
\setcounter{secnumdepth}{4}

\titleformat{\paragraph}
{\normalfont\normalsize\bfseries}{\theparagraph}{1em}{}
\titlespacing*{\paragraph}{0pt}{3.25ex plus 1ex minus .2ex}{1.5ex plus .2ex}

% Section styling
% ----------------------------------------------------------------------------------------------------------------------------------
\usepackage{sectsty}
\allsectionsfont{\color{dkblue}}

% Define source code highlighting
% ----------------------------------------------------------------------------------------------------------------------------------
\newcommand{\Hilight}{\makebox[0pt][l]{\color{cyan}\rule[-4pt]{0.65\linewidth}{14pt}}}

% Set the font to Helvetica
% ----------------------------------------------------------------------------------------------------------------------------------
\usepackage{helvet}
\renewcommand{\familydefault}{\sfdefault}

% Set margins
% ----------------------------------------------------------------------------------------------------------------------------------
\usepackage[top=.9in, bottom=1in, left=.5in, right=.5in]{geometry}

% Sections start a new page
% ----------------------------------------------------------------------------------------------------------------------------------
\let\stdsection\section
\renewcommand\section{\newpage\stdsection}

% Format paragraphs with no indent and a blank line between paragraphs
% ----------------------------------------------------------------------------------------------------------------------------------
\setlength\parindent{0pt}
\usepackage{parskip}

% Add page headers and footers
% ----------------------------------------------------------------------------------------------------------------------------------
\usepackage{fancyhdr}

\fancyhead[LE,RO]{\slshape \rightmark}
\fancyhead[LO,RE]{\slshape \leftmark}

\fancypagestyle{plain}
{
    \fancyhead{}
    \lhead[]{TABLE OF CONTENTS}
}

\lfoot[]{Crunchy Postgres Operator Quickstarts Guide\\DRAFT}
\cfoot[]{\ \\-\ \thepage\ -}
\rfoot[]{Crunchy Data Solutions, Inc.\\\today}
\pagestyle{fancy}

\renewcommand{\headrulewidth}{0.4pt}
\renewcommand{\footrulewidth}{0.4pt}

% Use framed package for admonition
% ----------------------------------------------------------------------------------------------------------------------------------
\usepackage{framed}

% ----------------------------------------------------------------------------------------------------------------------------------
% Begin document
% ----------------------------------------------------------------------------------------------------------------------------------
\begin{document}

% Create the title page
% ----------------------------------------------------------------------------------------------------------------------------------
\hypersetup{pageanchor=false}

\makeatletter
    \begin{titlepage}
        \begin{center}
            {\large \ }\\[18ex]
            {\huge \bfseries Crunchy Postgres Operator}\\[1ex]
            {\large Quickstarts Guide}\\[12ex]
            {\large DRAFT}\\[12ex]
            \includegraphics[width=6in]{/crunchy-docdynamo-postgres-operator/doc/output/latex/logo}\\[12ex]
            {\large Crunchy Data Solutions, Inc.}\\[1ex]
            {\large \today}
        \end{center}
    \end{titlepage}
\makeatother
\thispagestyle{empty}
\newpage

\hypersetup{pageanchor=true}

% Generate TOC
% ----------------------------------------------------------------------------------------------------------------------------------
\setcounter{tocdepth}{3}
\topskip0in
\thispagestyle{plain}
\renewcommand\contentsname{Table of Contents}
\tableofcontents

% ----------------------------------------------------------------------------------------------------------------------------------
% Content
% ----------------------------------------------------------------------------------------------------------------------------------

% 'Overview' Section
% ----------------------------------------------------------------------------------------------------------------------------------
\section{Overview}\label{/_overview}

There are currently \textbf{quickstart} scripts that seek to automate the deployment to popular Kube environments: * quickstart-for-gke.sh * quickstart-for-ocp.sh

The \textbf{quickstart-for-gke} script will deploy the operator to a GKE Kube cluster.

The \textbf{quickstart-for-ocp} script will deploy the operator to an Openshift Container Platform cluster.

Both scripts assume you have a StorageClass defined for persistence.

Pre-compiled versions of the Operator \textbf{pgo} client are provided for the x86\_64, Mac OSX, and Windows hosts.

% 'Quickstart' Section
% ----------------------------------------------------------------------------------------------------------------------------------
\section{Quickstart}\label{/_quickstart}

% 'Quickstart' Section, 'GKE/PKS' SubSection
% ----------------------------------------------------------------------------------------------------------------------------------
\subsection{GKE/PKS}\label{/_quickstart/_gke_pks}

The \textbf{quickstart-for-gke.sh} script will allow users to set up the Postgres Operator quickly on GKE including PKS. This script is tested on GKE but can be modified for use with other Kubernetes environments as well.

The script requires a few things in order to work:

\begin{itemize}
    \item wget utility installed 
    \item kubectl utility installed 
    \item StorageClass defined 
\end{itemize}
Executing the script will give you a default Operator deployment that assumes \textbf{dynamic} storage and a storage class named \textbf{standard}, things that GKE provides.

The script performs the following:

\begin{itemize}
    \item downloads the Operator configuration files 
    \item sets the \$HOME/.pgouser file to default settings 
    \item deploys the Operator Deployment 
    \item sets your .bashrc to include the Operator environment variables 
    \item sets your \$HOME/.bash\_completion file to be the \textbf{pgo} bash\_completion file 
\end{itemize}
% 'Quickstart' Section, 'Openshift Container Platform' SubSection
% ----------------------------------------------------------------------------------------------------------------------------------
\subsection{Openshift Container Platform}\label{/_quickstart/_openshift_container_platform}

A similar script for installing the operator on OCP is offered with similar features as the GKE script. This script is tested on OCP 3.7 with a StorageClass defined.

% ----------------------------------------------------------------------------------------------------------------------------------
% End document
% ----------------------------------------------------------------------------------------------------------------------------------
\end{document}
